\documentclass[11pt]{article}
\usepackage[margin=1in]{geometry}
\usepackage[T1]{fontenc}
\usepackage[utf8]{inputenc}
\usepackage{amsmath,amssymb,mathtools}
\usepackage{booktabs}
\usepackage{hyperref}
\usepackage{xcolor}
\usepackage{listings}

\hypersetup{
  colorlinks=true,
  linkcolor=blue!60!black,
  urlcolor=blue!60!black
}

\definecolor{codebg}{RGB}{248,248,248}
\definecolor{codeborder}{RGB}{210,210,210}
\definecolor{codecomment}{RGB}{80,120,80}
\definecolor{codekeyword}{RGB}{30,60,150}

\lstdefinestyle{reportcode}{
  backgroundcolor=\color{codebg},
  basicstyle=\ttfamily\small,
  breaklines=true,
  columns=fullflexible,
  commentstyle=\color{codecomment},
  frame=single,
  framerule=0.5pt,
  keywordstyle=\color{codekeyword},
  rulecolor=\color{codeborder},
  keepspaces=true,
  showstringspaces=false,
  tabsize=2
}

\lstset{style=reportcode, language=Python}

\setlength{\parindent}{0pt}
\setlength{\parskip}{0.5em}

\newcommand{\R}{\mathbb{R}}
\newcommand{\ysp}{y_{\mathrm{sp}}}
\newcommand{\uvec}{u}
\newcommand{\xvec}{x}
\newcommand{\Aaug}{A_{\mathrm{aug}}}
\newcommand{\Baug}{B_{\mathrm{aug}}}
\newcommand{\Caug}{C_{\mathrm{aug}}}
\newcommand{\Qout}{Q_{\mathrm{out}}}
\newcommand{\Rin}{R_{\mathrm{in}}}
\newcommand{\Hp}{H_p}
\newcommand{\Hc}{H_c}
\newcommand{\argmin}{\mathop{\mathrm{arg\,min}}}


\title{Polymer Online Re-Identification With Batch Ridge And RL Blend}
\date{2026-04-12}

\begin{document}
\maketitle

\section*{Experiment Surface}

The repository now exposes two polymer-only online re-identification
workflows:

\begin{itemize}
\item \texttt{RL\_assisted\_MPC\_reid\_batch\_unified.ipynb}: legacy comparison
      baseline
\item \texttt{RL\_assisted\_MPC\_reid\_batch\_v2\_unified.ipynb}: debug and
      improvement path
\end{itemize}

Both paths share the same runner stack:
\texttt{utils/reid\_batch.py},
\texttt{utils/reid\_batch\_runner.py},
\texttt{utils/plotting.py}, and
\texttt{utils/plotting\_core.py}.

\section*{Observer And Prediction Split}

Both legacy and v2 keep the observer nominal and fixed:
\[
A_{\mathrm{obs}} = A_0,\qquad
B_{\mathrm{obs}} = B_0,\qquad
C_{\mathrm{obs}} = C_0,\qquad
L = L_0.
\]

Only the MPC prediction model changes:
\[
A_{\mathrm{pred},k} = A_0 + \eta_k (A_{\mathrm{id},k} - A_0), \qquad
B_{\mathrm{pred},k} = B_0 + \eta_k (B_{\mathrm{id},k} - B_0).
\]

The RL policy acts only on the scalar blend factor $\eta_k \in [0,1]$.

\section*{Legacy Vs V2}

\subsection*{Legacy}

The legacy path keeps the original phase-1 scalar basis:

\begin{itemize}
\item one global physical-$A$ correction parameter
\item one physical-$B$ correction parameter per input column
\end{itemize}

Its blend-policy extras are raw:
\[
\eta_{k-1},\;
\|\varepsilon_k^{\mathrm{id}}\|_2,\;
\|\Delta A_k\|_F / \|A_0\|_F,\;
\|\Delta B_k\|_F / \|B_0\|_F,\;
\mathbb{1}_{\mathrm{id\_valid}}.
\]

\subsection*{V2}

The v2 path adds two richer basis families.

\paragraph*{Row/column basis}
\[
E_i = e_i e_i^\top A_0, \qquad
F_j = B_0 e_j e_j^\top.
\]

This yields one physical-state row basis for $A$ and one input-column basis
for $B$.

\paragraph*{Polymer block basis}

The physical states are partitioned into row blocks. For block $g$ with
row-mask projector $P_g$:
\[
E_g = P_g A_0.
\]

The $B$ basis remains one input-column basis per input.

V2 defaults to \texttt{basis\_family = rowcol}, while
\texttt{block\_polymer} remains available for ablation.

\section*{Candidate Guard Modes}

V2 replaces the previous implicit screening with explicit guard modes:

\begin{itemize}
\item \texttt{theta\_only}: enforce theta bounds only
\item \texttt{fro\_only}: enforce $\Delta A$ / $\Delta B$ Frobenius caps only
\item \texttt{both}: enforce both bound families
\end{itemize}

The default v2 choice is \texttt{fro\_only}. The legacy path keeps the more
restrictive effective behavior.

\section*{Blend-Policy State}

The base RL state remains the unified notebook state:
\[
s_k^{\mathrm{base}} = \texttt{build\_rl\_state}(\hat{x}_k, y_{\mathrm{sp},k}, u_k).
\]

Legacy appends the raw extras described above. V2 instead appends normalized
diagnostics:
\[
z_\eta = 2\eta_{k-1} - 1,
\]
\[
z_r = \mathrm{clip}\!\left(
\frac{\log(1 + \|\varepsilon_k^{\mathrm{id}}\|_2)}{s_r},
-c_r, c_r
\right),
\]
\[
z_A = \mathrm{clip}\!\left(\frac{\|\Delta A_k\|_F / \|A_0\|_F}{\Delta A_{\max}}, 0, 1.5\right),
\qquad
z_B = \mathrm{clip}\!\left(\frac{\|\Delta B_k\|_F / \|B_0\|_F}{\Delta B_{\max}}, 0, 1.5\right),
\]
plus normalized candidate-model ratios and the flags
\[
\mathbb{1}_{\mathrm{id\_valid}},\;
\mathbb{1}_{\mathrm{update\_success}},\;
\mathbb{1}_{\mathrm{fallback}}.
\]

\section*{Observer-Timing Ablation}

The observer remains nominal in every mode, but the timing of the correction is
now selectable in the re-identification runner:

\begin{itemize}
\item \texttt{legacy\_previous\_measurement}
\item \texttt{current\_measurement}
\end{itemize}

V2 defaults to \texttt{current\_measurement}. This changes only correction
timing, not observer matrices.

\section*{Runtime Loop}

\begin{enumerate}
\item build the nominal observer state and the current blend-policy state
\item choose the scalar RL action and map it to $\eta_k$
\item blend the nominal and identified physical models
\item solve MPC using the blended prediction model
\item step the plant
\item update the nominal observer
\item append the physical transition to the rolling identification buffer
\item every \texttt{id\_update\_period} samples, solve the batch
      identification problem on the most recent \texttt{id\_window} samples
\item validate the candidate according to the selected guard mode
\item keep the candidate if valid; otherwise retain the previous active model
\item store the RL transition and train the continuous agent
\end{enumerate}

\section*{Diagnostics}

The saved bundle now supports both legacy and v2 analysis. New v2 diagnostics
include:

\begin{itemize}
\item candidate vs active theta traces
\item candidate vs active $\Delta A$ / $\Delta B$ ratio traces
\item theta unclipped values
\item lower-bound and upper-bound hit masks
\item theta clipped fraction
\item candidate-valid flags
\item update-success and fallback markers
\item reward-vs-$\eta$ visualization
\end{itemize}

Legacy bundles remain plot-compatible because all new fields are optional.

\section*{Scope}

\begin{itemize}
\item Polymer only
\item Observer always nominal
\item The original \texttt{reid\_batch} workflow remains available unchanged
\item No changes to the matrix, structured-matrix, baseline, horizon, weights,
      residual, or combined methods
\end{itemize}

\end{document}
